% =============================================================================
\chapter{Ética, LGPD e limitações}\label{cap:etica}
% =============================================================================

Este capítulo consolida a análise de governança do projeto: a conformidade com a
LGPD \cite{brasil2018}, o respeito à licença da fonte, os vieses e limitações dos
dados e do método e os riscos de transformar análises estatísticas em decisões
automatizadas. A versão detalhada está em \code{ETHICS_LGPD.md}.

\section{Conformidade com a LGPD}\label{sec:lgpd}

Nos termos do art.~5º, inciso~I, da LGPD, dado pessoal é a ``informação
relacionada a pessoa natural identificada ou identificável''
\cite[art.~5º, I]{brasil2018}. A base analisada contém exclusivamente registros
públicos de execução no livro de ofertas à vista de BTC/USDT da Binance. O
\autoref{qua:lgpd} avalia cada campo.

\begin{quadro}[htbp]
\caption{\label{qua:lgpd}Avaliação dos campos da base quanto à presença de dados pessoais}
\begin{corpotabela}
\begin{tabularx}{\textwidth}{|L{3.0cm}|Y|L{4.6cm}|}
\hline
\textbf{Campo} & \textbf{Natureza} & \textbf{Contém dado pessoal?} \\ \hline
\code{trade_id} & Número sequencial atribuído pelo motor de casamento da corretora & Não. Identifica um evento de mercado, sem vínculo público com conta, CPF, CNPJ ou carteira \\ \hline
\code{exchange_ts} (\code{ts}) & Instante da execução pública, em UTC & Não. Metadado temporal de mercado \\ \hline
\code{price}, \code{quantity} & Preço e quantidade executados no livro público & Não. Informação pública de formação de preço \\ \hline
\code{aggressor_side} (\code{side}) & Lado do agressor (compra ou venda) & Não. Dinâmica impessoal de oferta e demanda \\ \hline
\code{exchange}, \code{symbol} & Constantes técnicas (\code{binance}, \code{BTCUSDT}) & Não \\ \hline
\code{ingest_ts} & Carimbo de coleta do fornecedor, 100\,\% inválido & Não; além disso, foi descartado na limpeza \\ \hline
\end{tabularx}
\end{corpotabela}
\fonteautores[Elaborado]
\end{quadro}

Não há titular identificável nem dado pessoal sensível (art.~5º, II). Como a LGPD
regula o tratamento de dados pessoais, o tratamento de dados de mercado anônimos
fica fora do seu escopo. Mesmo que houvesse dado pessoal, o art.~4º, inciso~II,
alínea~b, afasta a aplicação da lei ao tratamento realizado para fins
exclusivamente acadêmicos, ressalvados os arts.~7º e 11.

Ainda assim, o projeto aplica, por boa prática, os princípios do art.~6º. Pelo
princípio da \textbf{necessidade}, foram removidas as colunas constantes,
redundantes e corrompidas, mantendo-se apenas o necessário às perguntas de
negócio. Pelo princípio da \textbf{qualidade dos dados}, a base passou por
contrato de qualidade e por validação cruzada com a fonte primária. Pelo
princípio da \textbf{transparência}, todas as decisões de tratamento estão
documentadas.

\section{Risco residual de reidentificação}\label{sec:reidentificacao}

Embora a base isolada seja impessoal, em uma mesa real existiria um risco
conceitual de reidentificação indireta se os registros públicos fossem cruzados
com logs internos de ordens da própria mesa, com cadastros de clientes
(processos de identificação de clientes, ou KYC) ou com movimentações em
\textit{blockchain} próximas ao instante de uma ordem. O projeto mitiga esse risco
de três formas: opera exclusivamente sobre o fluxo público, sem nenhum
cruzamento com cadastros, carteiras ou logs de clientes; trabalha
predominantemente com agregações temporais (barras de 1 minuto e 1 segundo); e
reconstrói ordens agressoras apenas por contiguidade temporal, sem atribuição de
autoria.

\section{Licença e reprodutibilidade}\label{sec:licenca}

A base está sob licença de uso pessoal e acadêmico, com redistribuição proibida
\cite{glitchlist2026}. Os controles técnicos adotados são:

\begin{alineas}
  \item \textbf{sem redistribuição}: o arquivo \code{.gitignore} bloqueia toda a
        pasta \code{data/}; nenhum Parquet bruto ou processado é versionado;
  \item \textbf{download na origem e imutabilidade}: o código fixa o
        \textit{commit} da base no Hugging Face e confere o SHA-256 antes de
        qualquer processamento;
  \item \textbf{finalidade acadêmica}: código, documentação e painel destinam-se
        à avaliação da unidade curricular e à demonstração metodológica em um
        cenário simulado;
  \item \textbf{fonte primária apenas para validação}: os candles da Binance
        \cite{binance2025} são usados somente para reconciliar a base.
\end{alineas}

As figuras deste documento mostram apenas agregações e estatísticas derivadas,
não os registros da base.

\section{Vieses e limitações}\label{sec:vieses}

O recorte escolhido traz vieses que limitam a generalização dos resultados. O
\autoref{qua:vieses} resume cada um e o seu impacto na interpretação.

\begin{quadro}[htbp]
\caption{\label{qua:vieses}Vieses e limitações dos dados e do método}
\begin{corpotabela}
\begin{tabularx}{\textwidth}{|L{3.3cm}|Y|L{5.0cm}|}
\hline
\textbf{Dimensão} & \textbf{Limitação} & \textbf{Impacto na interpretação} \\ \hline
Seleção temporal & Apenas 7 dias (19 a 25/02/2025), marcados pelo ataque à Bybit e por queda de 13,5\,\%; saldo vendedor em todos os dias & Os padrões não representam semanas neutras ou de alta \\ \hline
Única corretora & Somente a Binance \textit{spot}, a de maior liquidez global & Outras corretoras e mercados de derivativos podem ter horários de pico e profundidade diferentes \\ \hline
Único par & Somente BTC/USDT, o par mais líquido do mercado & Outros criptoativos e pares contra moeda fiduciária têm \textit{spreads} e lacunas de liquidez maiores \\ \hline
Amostra horária curta & Apenas 7 observações por hora do dia, das quais 3 em estresse & As medianas horárias são indicativas, não definitivas \\ \hline
Reconstrução de ordens & Negócios do mesmo milissegundo e lado são fundidos em uma ordem & Ordens independentes simultâneas podem ser somadas \\ \hline
Qualidade do fornecedor & \code{ingest_ts} 100\,\% inválido e período declarado incorreto & Impede qualquer análise de latência de coleta; exigiu validação cruzada integral \\ \hline
Sem livro de ofertas & A base contém apenas negócios, sem cotações nem profundidade & Não é possível medir o \textit{spread} cotado ou a liquidez em repouso, apenas o impacto realizado \\ \hline
\end{tabularx}
\end{corpotabela}
\fonteautores[Elaborado]
\end{quadro}

\section{Riscos do uso de sinais de mercado}\label{sec:riscos}

Mesmo com uma base limpa e validada, transformar análises em decisões
automatizadas envolve riscos que a mesa precisa conhecer:

\begin{alineas}
  \item \textbf{ilusão preditiva diante dos custos}: acurácia ou ROC-AUC acima de
        50\,\% não implica lucro; o ganho bruto dos sinais direcionais é muito menor
        que as taxas, e operá-los geraria prejuízo por excesso de giro;
  \item \textbf{quebra de regime}: modelos treinados em dias calmos subestimam a
        escala da volatilidade e o tamanho das ordens durante choques; a
        reavaliação contínua em janelas \textit{walk-forward} e limites rígidos
        baseados em volatilidade são indispensáveis;
  \item \textbf{vazamento de futuro e mineração de dados}: interpolação de lacunas,
        normalização antes da divisão entre treino e teste ou atributos com
        \code{shift(-1)} criam acurácia artificial; o projeto evita esses erros e os
        testa automaticamente;
  \item \textbf{impacto próprio e retroalimentação}: uma mesa que executa ordens
        grandes sempre no mesmo horário, ou que reage mecanicamente ao OFI, passa
        a gerar ela própria o desequilíbrio observado e aumenta o próprio
        \textit{slippage};
  \item \textbf{uso responsável}: os modelos são recomendados exclusivamente como
        \textbf{apoio à decisão}, com supervisão humana, para dimensionar janelas
        de execução e calibrar limites, e não como motor autônomo de
        negociação.
\end{alineas}
